%!TEX root = ../main.tex
\chapter{付録}

\section{見かけの位置の導出}\label{sec:appendix-apparent-depth}

カメラ座標系においてカメラ位置を$z$軸上$\left(0, 0, H\right)$に、また水面を平坦な$xy$平面(法線ベクトルが一定、すなわち静水面とする)に配置する。
水中に存在する各三次元Gaussianプリミティブの位置$\bm{p}$は$\left(x, y, z\right)$(ただし$z<0$)で表される。
簡単化のため、問題を2次元化し、位置を$r = \sqrt{x^2 + y^2}$として$rz$-平面上に射影する(\cref{fig:rz}を参照)。

幾何光学原理およびスネルの法則(Snell's law)より、水中点$\bm{p}$を観測するカメラ視点から見た「見かけの位置」は、観測方向に依存せず定まる\cite{nassar1994_ApparentDepth,Missailidis2025_apparentDepth-leading}。
ここではその詳細な導出を示す。

Gaussian中心の見かけの位置を$\bm{p}'(r', z')$とする。
カメラ位置と見かけの位置$\bm{p}'$を結ぶ入射線と水面との交点を$(s, 0)$($0 \leq s < r$)とおく。
スネルの法則(空気屈折率$1$、水屈折率$n$、空中角$\theta_a$、水中角$\theta_w$)は次式で表される:
\begin{equation}\label{eq:-snell}
  \sin \theta_a = n \sin \theta_w
\end{equation}
上記関係から、$s$に関する四次方程式が導かれる:
\begin{equation}\label{eq:-quartic-s}
\begin{split}
  (1-n^2)s^4 &+ 2(n^2-1)rs^3 \\
  &+ ((1-n^2)r^2 - n^2H^2 + z^2)s^2 \\
  &+ 2n^2H^2rs - n^2H^2r^2 = 0
\end{split}
\end{equation}
ニュートン法やFerrariの公式で、$0 \leq s < r$の範囲で解くことで、光線の水面との交点位置$s$が求められる。

わずかなずれ$\Delta s$を仮定し、$s+\Delta s$の交点、および$\theta_w+\Delta \theta_w$の微小変化を考える。
辺$IP'$と鉛直･水平線がつくる三角形の関係より:
\begin{equation}\label{eq:-triangle-IP'C'}
   r' - s = -z' \tan \theta_a 
\end{equation}
また、$\triangle I'P'C'$より:
\begin{align}\label{eq:-triangle-I'P'C'}
   r' - ( s + \Delta s ) & = -z' \tan \left(\theta_a + \Delta \theta_a \right)  \notag \\
   \iff (r' -s - \Delta s) \left( \cos \theta_a - \Delta \theta_a \sin \theta_a \right) & = -z' \left( \sin \theta_a + \Delta \theta_a \cos \theta_a  \right) 
\end{align}
\cref{eq:-triangle-IP'C'}と\cref{eq:-triangle-I'P'C'}から$z'$を消去し、$\Delta \to 0$の極限をとると:
\begin{equation}\label{eq:-limit-r'}
  r' = s - \sin \theta_a \cos \theta_a \frac{d s}{d \theta_a}
\end{equation}
が得られる。s
これを再び\cref{eq:-triangle-IP'C'}に用い、$z'$で表すと:
\begin{equation}\label{eq:-z'_ds/dtheta_i}
  z' = \frac{d s}{d \theta_a} {\cos \theta_a}^2
\end{equation}
さらに、三角形$\triangle IPC$から:
\begin{equation}\label{eq:-r-s}
  r - s = - z \tan \theta_w
\end{equation}
これを$\theta_w$で微分すると:
\begin{equation}\label{eq:-ds/dtheta_r}
  \frac{d s}{d \theta_w} = z \cdot \frac{1}{{\cos \theta_w}^2}
\end{equation}
また、スネルの法則\eqref{eq:-snell}を$\theta_a$で微分すると:
\begin{equation}\label{eq:-dtheta_r/dtheta_i}
  \frac{d \theta_w}{d \theta_a} = \frac{1}{n} \cdot \frac{\cos \theta_a}{\cos \theta_w}
\end{equation}
\eqref{eq:-dtheta_r/dtheta_i},\eqref{eq:-ds/dtheta_r}より:
\begin{equation}\label{eq:-ds/dtheta_i}
  \frac{d s}{d \theta_a} = \frac{z}{n} \cdot \frac{\cos \theta_a}{{\cos \theta_w}^3}
\end{equation}
以上の式を整理すると、見かけの位置$(r', z')$はそれぞれ:
\begin{subequations}\label{eq:apparent-position}
  \begin{empheq}[left=\empheqlbrace]{align}
    r' &= r + (n^2 -1) \cdot z \cdot \tan^3 \theta_w \label{eq:r'} \\
    z' &= \frac{1}{n} \cdot \frac{{\cos \theta_a}^3}{{\cos \theta_w}^3} \cdot z \label{eq:z'}
\end{empheq}
\end{subequations}
と表される。

